\section{Conclusion}
This paper develops a unified, privacy-preserving framework that advances anti–money laundering (AML) detection from isolated, rule-based monitoring toward coordinated, explainable, and action-oriented financial crime prevention. By integrating graph-based federated learning, cross-bank Personalized PageRank, interpretable group identification, and hierarchical reinforcement learning, the proposed system addresses four long-standing challenges in AML: institutional data silos, lack of interpretability, weak decision support following detection, and severe class imbalance.

Our federated graph learning framework enables institutions to learn global laundering behaviors without sharing sensitive data, significantly reducing both false alarms and missed detections. Empirical analysis demonstrates that federated training improves performance across countries, while graph-based relational modeling captures complex behavioral dependencies that feature-driven methods overlook. Through cross-bank PPR, the system transforms raw detection signals into interpretable group structures, uncovering coordinated laundering patterns such as fan-out, gather–scatter, and hybrid schemes distributed across institutions and borders. These interpretable clusters strengthen investigative transparency and provide regulators with clearer evidence of network-level criminal activity.

Furthermore, the proposed hierarchical reinforcement learning mechanism shifts AML from static alerting to adaptive intervention. By optimizing when to escalate, monitor, or intervene, the framework aligns detection outputs with operational constraints, regulatory requirements, and cost considerations. This represents a meaningful departure from traditional rule-based escalation policies, enabling timely and economically efficient responses to suspicious activity.

Our results show that combining privacy-preserving collaboration with structure-aware learning and principled decision-making yields substantial gains in detection accuracy, interpretability, and operational effectiveness. The insights also highlight broader managerial implications: institutions that can securely leverage cross-organizational signals and act adaptively will be better positioned to counter increasingly sophisticated laundering threats. Future research may explore extending this framework to real-time streaming environments~(\cite{zheng2024online, zheng2024mulan}), adversarial laundering strategies, and multi-modal datasets incorporating customer profiles, text reports, and cross-border regulatory information. By bridging detection, explanation, and action, this work contributes a new foundation for scalable, transparent, and collaborative financial crime prevention in a globally interconnected economy.
